\chapter{Testing}
Testing using pytest
\begin{minted}{python}
    import pytest
    import EmailAddressValidator
    
    
    class TestSportradarEmailAddressValidator:
        @pytest.mark.parametrize(
            "input_email, expected_validity",
            [
                ("someone@sportradar.com", True),
                ("anotherone@sport.at", False),
                ("dead.red", False),
            ]
        )
        def test_get_validity(self, input_email, expected_validity):
            email_validator = SportRadarEmailAddressValidator(input_email)
            assert email_validator.get_validity() == expected_validity
\end{minted}

\section{assert}
Assert should be interpreted as \textit{to make sure}

\section{fixures}
    
    \paragraph{autouse=True}
        \begin{minted}{python}
            # tests/conftest.py
            import pathlib
            import pytest
            
            @pytest.fixture(autouse=True)
            def chdir_to_test_dir(request, monkeypatch):
                monkeypatch.chdir(pathlib.Path(request.fspath).parent)
        \end{minted}
    \paragraph{With yieldd}
        \begin{minted}{python}
            @pytest.fixture
            def conf_fixture():
                # load config file
                yield # test runs here
                # remove the config structure
        \end{minted}
        
    
\section{conftest.py}
    A file that runs before the tests from the folder in which it is to deeper.
    The followng code adds each test's file path to tests directory for example if a test file is in tests/storage/type/test\_file.py, it changes the working directory from tests/test\_file.py to tests/storage/type/test\_file.py
    \begin{minted}{python}
        # tests/conftest.py
        import pathlib
        import pytest
        
        @pytest.fixture(autouse=True)
        def chdir_to_test_dir(request, monkeypatch):
            monkeypatch.chdir(pathlib.Path(request.fspath).parent)
    \end{minted}